\documentclass[
  12pt,
  oneside,
  a4paper,
  chapter=TITLE,
  section=TITLE,
  english,
  brazilian
]{abntex2}

\usepackage{iftex}
\ifPDFTeX
  \usepackage[T1]{fontenc}
  \usepackage[utf8]{inputenc}
  \usepackage{mathptmx}
  \renewcommand{\sfdefault}{ptm}
  \renewcommand{\ttdefault}{ptm}
\else
  \usepackage{fontspec}
  \IfFontExistsTF{Times New Roman}{%
    \setmainfont{Times New Roman}\setsansfont{Times New Roman}\setmonofont{Times New Roman}}{%
    \IfFontExistsTF{Liberation Serif}{%
      \setmainfont{Liberation Serif}\setsansfont{Liberation Serif}\setmonofont{Liberation Serif}}{%
      \setmainfont{TeX Gyre Termes}\setsansfont{TeX Gyre Termes}\setmonofont{TeX Gyre Termes}}}
\fi
\usepackage{indentfirst}
\usepackage{graphicx}
\usepackage{xcolor}
\usepackage{tabularx}
\usepackage{booktabs}
\usepackage{placeins}
\usepackage{microtype}
\usepackage[alf,abnt-etal-cite=3,abnt-etal-text=it]{abntex2cite}
\usepackage{url}
\urlstyle{same}
\Urlmuskip=0mu plus 1mu
% NBR 6023:2018: endereço eletrônico sem os sinais < >
\def\UrlLeft{}\def\UrlRight{}

\let\oldabntrefinfo\abntrefinfo
\renewcommand{\abntrefinfo}[3]{\oldabntrefinfo{#1}{#1}{#3}}

\renewcommand{\arraystretch}{1.15}
\newcolumntype{Y}{>{\raggedright\arraybackslash}X}
\newcolumntype{P}[1]{>{\raggedright\arraybackslash}p{#1}}
\setlength{\emergencystretch}{3em}
\raggedbottom
\renewcommand{\topfraction}{0.9}
\renewcommand{\bottomfraction}{0.9}
\renewcommand{\textfraction}{0.08}
\renewcommand{\floatpagefraction}{0.75}
\setcounter{topnumber}{3}
\setcounter{bottomnumber}{3}
\setcounter{totalnumber}{6}
\makeatletter
\setlength{\@fptop}{0pt}
\makeatother
\AtBeginDocument{%
  \renewcommand{\fonte}[1]{\par\vspace{3pt}{\footnotesize\raggedright\noindent Fonte: #1\par}\vspace{4pt}}%
  \renewcommand{\nota}[1]{\par\vspace{2pt}{\footnotesize\raggedright\noindent Nota: #1\par}\vspace{4pt}}%
}

% Quadro como flutuante próprio (mesmo padrão da E1)
\newcommand{\quadroname}{Quadro}
\newcommand{\listofquadrosname}{Lista de quadros}
\newfloat[chapter]{quadro}{loq}{\quadroname}
\newlistof{listofquadros}{loq}{\listofquadrosname}
\newlistentry{quadro}{loq}{0}
\setfloatadjustment{quadro}{\centering}
\counterwithout{quadro}{chapter}
\renewcommand{\cftquadroname}{\quadroname\space}
\renewcommand*{\cftquadroaftersnum}{\hfill\textendash\hfill}

\newcommand{\preencher}[1]{\textcolor{red}{[#1]}}

% Código da tela principal (errata da E2: opcional, não entregue nesta etapa)
\newif\ifcodigo
\codigofalse

% Imagem com fallback: procura em ./figuras/, depois nas figuras da E1
\newcommand{\efig}{../../entrega1/relatorio/figuras}
\newcommand{\img}[2]{%
  \IfFileExists{figuras/#2}{\includegraphics[#1]{figuras/#2}}{%
    \IfFileExists{\efig/#2}{\includegraphics[#1]{\efig/#2}}{%
      \fbox{\parbox[c][3.5cm][c]{0.9\linewidth}{\centering
        \preencher{inserir figuras/\detokenize{#2}}}}}}}

\newcommand{\fig}[5][0.95]{%
  \begin{figure}[!htbp]
    \centering
    \caption{#4}\label{fig:#2}
    \img{width=#1\textwidth,height=0.42\textheight,keepaspectratio}{#3}
    \fonte{#5}
  \end{figure}\FloatBarrier}
\newcommand{\capt}[1]{Secretaria de Economia do DF (#1), acesso em 18 out. 2026}
\newcommand{\autores}{Elaborado pelos autores (2026).}

% Antes e depois lado a lado, dentro do quadro de decisão
\newcommand{\antesdepois}[4]{%
  \noindent
  \begin{minipage}[b]{0.48\textwidth}\centering
    \img{width=\linewidth,height=6.5cm,keepaspectratio}{#1}\par
    {\footnotesize Antes: #2}
  \end{minipage}\hfill
  \begin{minipage}[b]{0.48\textwidth}\centering
    \img{width=\linewidth,height=6.5cm,keepaspectratio}{#3}\par
    {\footnotesize Depois: #4}
  \end{minipage}\par\vspace{6pt}}

% Corpo QOC do quadro de decisão
\newcommand{\qoc}[6]{%
  \begin{tabularx}{\textwidth}{|P{3.6cm}|Y|}\hline
    Questão (Q)                 & #1 \\\hline
    Opção adotada (O)           & #2 \\\hline
    Alternativa descartada (O)  & #3 \\\hline
    Requisito (C)               & #4 \\\hline
    Origem (C)                  & #5 \\\hline
    Motivo                      & #6 \\\hline
  \end{tabularx}}

% ---------------------------------------------------------------------
% Dados do trabalho
% ---------------------------------------------------------------------
\titulo{Protótipo do fluxo de solicitação de emissão da guia de ITBI no portal da Secretaria de Economia do DF}
\autor{%
  Ícaro Marra Souza\\
  Karoline Rodrigues Costa\\
  Leonardo Brito Gomes\\
  Luiz Gustavo Silva Carvalho}
\local{Brasília}
\data{nov. 2026}
\preambulo{Memorial referente à entrega 2 do projeto da Unidade
Curricular de Engenharia de Usabilidade, apresentado à Universidade do Distrito
Federal como requisito de avaliação, sob orientação da Profa. Dra. Carolina
Gonçalves Abreu.}

\definecolor{blue}{RGB}{41,5,195}
\makeatletter
\hypersetup{
  pdftitle={\@title},
  pdfauthor={\@author},
  pdfsubject={UCEngUs -- entrega 2},
  pdfcreator={LaTeX com abnTeX2},
  pdfkeywords={usabilidade}{acessibilidade}{protótipo}{design de interface},
  colorlinks=true,
  linkcolor=blue,
  citecolor=blue,
  filecolor=magenta,
  urlcolor=blue,
  bookmarksdepth=4}
\makeatother

\renewcommand{\imprimircapa}{%
  \begin{capa}%
    \center
    \IfFileExists{../../assets/marca/undf-vertical.pdf}{%
      \includegraphics[width=4.5cm]{../../assets/marca/undf-vertical.pdf}\par}{%
      \IfFileExists{undf-vertical.pdf}{\includegraphics[width=4.5cm]{undf-vertical.pdf}\par}{%
        \textcolor{red}{[marca UnDF vertical, mínimo 4 cm]}\par}}
    \vspace*{1cm}
    \ABNTEXchapterfont\large\imprimirautor

    \vfill
    \begin{center}
    \ABNTEXchapterfont\bfseries\LARGE\imprimirtitulo
    \end{center}
    \vfill

    \large\imprimirlocal

    \large\imprimirdata

    \vspace*{1cm}
  \end{capa}}

\setlength{\parindent}{1.25cm}
\setlength{\parskip}{0.2cm}

\makeindex

% =====================================================================
\begin{document}
\selectlanguage{brazil}
\frenchspacing

\pretextual

\imprimircapa
\imprimirfolhaderosto

\begin{resumo}
Este trabalho descreve, sobretudo, a proposta de redesenho de interface para o fluxo
crítico que vai da página inicial de serviços da
Secretaria de Estado de Economia do Distrito Federal (SEEC) até a abertura do
formulário de Atendimento Virtual para solicitar a emissão da guia do Imposto
sobre a Transmissão de Bens Imóveis (ITBI). O processo partiu, assim, de esboços feitos
à mão em sala, passou pelo mapa de navegação e pelo wireframe da tela principal
e terminou em um protótipo navegável no Figma; sendo que foram desenhadas telas responsivas, as quais podem ser acessadas tanto no celular quanto no computador. No caso, o fluxo contempla o caminho principal pela busca ou pelos
serviços mais procurados, uma rota alternativa por assunto e os estados de
validação, carregamento, vazio, erro e serviço fora do ar. Dado que o serviço
atual atravessa três sites diferentes e gera confusão de vocabulário e de
identidade visual, optamos por adotar os tokens e componentes do Padrão Digital
de Governo (DSGov), com adaptações pontuais para o Distrito Federal. As
cinco decisões de design registradas no formato QOC unem cada escolha aos
achados, às sessões e aos vinte requisitos da entrega 1, de modo cobrir a estrutura
da tela inicial, o acesso precário em celular, a linguagem simples voltada à
tarefa, a unificação visual com contraste acessível e a antecipação de login e
documentos antes de um formulário já pré-preenchido.

Palavras-chave: usabilidade, acessibilidade, protótipo, design de interface,
serviço público digital.
\end{resumo}

\listoffigures*
\cleardoublepage
\listofquadros*
\cleardoublepage
\ifcodigo
\listoftables*
\cleardoublepage
\fi

\begin{siglas}
  \item[DSGov]  Design System do Governo Federal
  \item[HTA]    \textit{Hierarchical Task Analysis} (Análise Hierárquica de Tarefas)
  \item[ITBI]   Imposto sobre a Transmissão de Bens Imóveis \textit{inter vivos}
  \item[LBI]    Lei Brasileira de Inclusão da Pessoa com Deficiência
  \item[NVDA]   \textit{NonVisual Desktop Access}
  \item[QOC]    \textit{Questions, Options and Criteria} (Questões, opções e critérios)
  \item[RA]     Requisito de acessibilidade
  \item[RU]     Requisito de usabilidade
  \item[SEEC]   Secretaria de Estado de Economia do Distrito Federal
  \item[UnDF]   Universidade do Distrito Federal
  \item[WCAG]   \textit{Web Content Accessibility Guidelines}
\end{siglas}

\tableofcontents*
\cleardoublepage

\textual

% =====================================================================
\chapter{Introdução}\label{cap:introducao}
% =====================================================================

Aborda-se, neste trabalho, uma proposta de definição da interface para o
fluxo crítico analisado na entrega 1 (E1), no qual o cidadão precisa encontrar
e abrir o formulário de solicitação da guia de ITBI no portal da Secretaria de
Estado de Economia do Distrito Federal \cite{seec}. Assim, seguimos a mesma linha da
ISO 9241-11:2018 \cite{iso9241-11-2018} adotada no relatório anterior, e também da perspectiva iterativa de design centrado no humano da ISO
9241-210:2019 \cite{iso9241-210-2019}, em que a solução de interface surge da compreensão do contexto de uso e dos requisitos especificados.
Por isso, o critério que guiou esta entrega foi a rastreabilidade, sendo que as
decisão tomada no protótipo se liga a pelo menos um requisito de
usabilidade (RU01 a RU08) ou de acessibilidade (RA01 a RA12) e ao achado (A1 a
A12) ou à sessão (S1 a S5) que originou o problema.

No diagnóstico da E1, ficou bem claro que a dificuldade das pessoas não vinha de um
único detalhe estético, sendo que o principal gargalo estava na descoberta da
rota e no vocabulário usado pelo órgão. Para chegar ao formulário hoje, o
usuário sai do portal institucional (\texttt{economia.df.gov.br}), passa pelo
catálogo da Receita (\texttt{receita.fazenda.df.gov.br}) e depois pelo
Atendimento Virtual (\texttt{atendimentovirtual.fazenda.df.gov.br}), passando
por três identidades visuais distintas e por rótulos ligados ao ato jurídico em
vez da tarefa real. Nas cinco sessões da E1, só uma pessoa conseguiu chegar ao
formulário sem ajuda, o que reforça que encurtar o caminho e traduzir os
rótulos era um ponto muito importante.

O Quadro~\ref{qua:ficha} traz os aspectos levantados na entrega 1 bem como os dados de
acesso ao protótipo navegável.

\begin{quadro}[!htbp]
  \caption{Dados provenientes da entrega 1 e link do protótipo}\label{qua:ficha}
  \centering\small
  \begin{tabularx}{\textwidth}{|P{3.6cm}|Y|}\hline
    Fluxo crítico & Encontrar e abrir, a partir da entrada de serviços da Secretaria de Economia do DF, o formulário do Atendimento Virtual para solicitar a emissão da guia de ITBI. \\\hline
    Início e término & Início na página de busca de serviços da Receita do DF, sem login prévio, sendo um equivalente redesenhado da entrada por \texttt{economia.df.gov.br}. Término na exibição do formulário ``Pedir a guia do ITBI'', já com Assunto e Tipo de atendimento selecionados. Como complemento demonstrativo no protótipo, incluíram-se também o estado de validação de campos e a confirmação de pedido registrado. \\\hline
    Etapas da HTA & 1 Acessar o site da Secretaria (1.1, 1.2, 1.3); 2 Acessar o Atendimento Virtual (2.1 e variante 2.1'); 3 Selecionar o serviço ITBI (3.1). \\\hline
    Classes e personas & C1 familiarizado com a web (Personas 3 Camila e 7 Rafael); C2 eventual com baixo letramento digital (Personas 1 Antônio e 4 Rosimeire); C3 acesso ou percepção restrita (Personas 2 Júlia, 5 Valdir e 6 Edilene). \\\hline
    Ferramenta e acesso & \url{https://www.figma.com/proto/ObXREG7SsNPGVHg9CIqmD4/UCEngUs-E2-GrupoB---ITBI--DSGov-?node-id=3-8&t=yMMosWHD1RiHZDEi-1}. \\\hline
  \end{tabularx}
  \fonte{\autores}
\end{quadro}\FloatBarrier

% =====================================================================
\chapter{Esboços da tela principal}\label{cap:esbocos}
% =====================================================================

A tela principal escolhida é a Tela 01 (Busca de serviços), pois ela ela concentra o ponto mais crítico do fluxo porque é logo na chegada que o
cidadão tenta traduzir uma necessidade concreta, que é pagar o imposto de um
imóvel recém-comprado, em algum caminho reconhecível na interface. Na E1,
quatro das cinco pessoas falharam exatamente nessa transição entre entrada e
descoberta do serviço, pois S1 precisou entrar e sair de menus até achar um link no
meio do texto, S2 e S4 se perderam em itens vizinhos do catálogo, S3 desistiu e S5 tentou a barra de busca cinco vezes sem
sucesso. Além do mais, a Tela 01 possui o campo de
entrada principal com regras de validação (campo vazio e texto curto), servindo
de base para o wireframe (capítulo~\ref{cap:wireframe}) e para a adaptação
entre celular e computador.

A partir da exploração dos protótipos de baixa fidelidade feitos a mão, consolidamos duas alternativas bem distintas para a tela
principal, sendo a Alternativa A 
(Figura~\ref{fig:esboco-a}), e a B (Figura~\ref{fig:esboco-b}).

\fig[0.78]{esboco-a}{esboco_karoline.jpg}{Esboço da alternativa A da tela principal feita pela Karoline Costa}{\autores}
\fig[0.78]{esboco-b}{esboco_leonardo.jpg}{Esboço da alternativa B da tela principal feira pelo Leonardo Brito}{\autores}

Quanto à concepção, a alternativa A foca em uma tela de entrada única e
direta, pois logo abaixo do cabeçalho, a barra de busca aparece em destaque com
campo visível e lupa, seguida imediatamente por uma lista vertical curta com as
tarefas mais procuradas. No caso, a ideia é diminuir o golfo de execução
\cite{norman2013}, colocando o objetivo do usuário logo no início da tela, como
o próprio participante de S3 sugeriu na entrevista, e ainda oferecer uma rota
secundária por assunto para quem prefere navegar clicando.

Já a alternativa B mantém uma arquitetura mais próxima da lógica hierárquica
do portal atual, organizando a entrada por menu retrátil e agrupamentos por
perfil de público (\textit{Atendimento Cidadão}, \textit{Atendimento Empresa})
e categorias de tributos (\textit{Serviços} $\rightarrow$ \textit{Imóveis}
$\rightarrow$ \textit{ITBI}), sem propociar enfoque à busca. Apesar de parecer
familiar a quem já conhece a divisão interna da Secretaria, essa estrutura
reproduz dois problemas medidos na E1, sendo que força escolhas sequenciais em múltiplos
níveis, aumentando o tempo de decisão pela Lei de Hick-Hyman \cite{hick1952}, e
exige que a pessoa saiba de antemão em qual opção a emissão do imposto foi
colocado.

% =====================================================================
\chapter{Fluxo de navegação}\label{cap:navegacao}
% =====================================================================

Com a estrutura da tela inicial definida, desenvolvemos o mapa de navegação
do início até a abertura do formulário (Figura~\ref{fig:fluxo-nav}). No serviço
real analisado na E1, o caminho sem erro requeria cinco cliques e quatro
carregamentos de página espalhados em três subdomínios, sendo que a única pessoa
que concluiu (S1) precisou de sete cliques pela variante do catálogo. No
redesenho, as etapas 1.1 a 1.3 da HTA (menus em cascata e escolha de perfil) e
a etapa 3.1 (escolha do serviço) foram trazidas para a porta de entrada em um
único domínio. Assim, a pessoa passa a contar com três caminhos paralelos para
o mesmo destino, pois ela pode pesquisar por palavra-chave (Tela 01 $\rightarrow$ Tela 02),
clicar no atalho de ``Mais procurados'' (Tela 01 $\rightarrow$ Tela 04)
ou navegar por assunto quando não quiser digitar (Tela 01 $\rightarrow$ Tela 03
$\rightarrow$ Tela 04).

Outro ponto importante foi reorganizar a passagem até o formulário, dado que o
diagnóstico mostrou que omitir a exigência de login e de documentos até o
final do percurso gera frustração e abandono tardio (achado A8). Por isso, o fluxo redesenhado posiciona a autenticação única pelo gov.br
(Tela 04) e uma tela intermediária de orientação prévia (Tela 05) logo antes do formulário. Nessa etapa, o cidadão já fica sabendo que
vai precisar do número da matrícula do imóvel, do contrato ou escritura e do
valor da compra.

Além do caminho de melhor caso, mapeamos os estados de erro, validação, carregamento,
vazio e indisponibilidade previstos para uma navegação realista, como a busca com
campo vazio (Tela 01a), busca com menos de três letras (Tela 01b), carregamento
da consulta (Tela 01c), busca fora do ar mantendo a lista de atalhos clicável
(Tela 01d), busca sem resultado (Tela 02a), carregamento e cancelamento ou falha
do login no gov.br (Telas 04a e 04b) e formulário indisponível (Tela 06a).

Quanto ao encerramento, o fluxo crítico
delimitado na E1 se encerra com a na exibição do formulário aberto com o serviço de ITBI
já escolhido (Tela 06), como mostra o diagrama da
Figura~\ref{fig:fluxo-nav}. Ainda assim, incluímos no protótipo navegável do
Figma as telas 06b (validação de campos com erro no formulário) e
06c (confirmação de pedido registrado com protocolo), de modo a funcionar como
demonstração complementar de prevenção e recuperação de erro. O Quadro~\ref{qua:telas} mostra, sobretudo, as
telas, sua correspondência na HTA e o papel de cada uma no fluxo.

\begin{figure}[!p]
  \centering
  \caption{Fluxo de navegação do protótipo}\label{fig:fluxo-nav}
  \img{width=\textwidth,height=0.8\textheight,keepaspectratio}{fluxo_navegacao.png}
  \fonte{\autores}
\end{figure}

\begin{quadro}[!htbp]
  \caption{Telas do protótipo, etapa da HTA e comportamento previsto}\label{qua:telas}
  \centering\footnotesize
  \begin{tabularx}{\textwidth}{|P{1.4cm}|P{3.6cm}|P{2.3cm}|Y|}\hline
    Tela & Nome & Etapa da HTA & Tipo e comportamento no fluxo \\\hline
    Tela 01 & Busca de serviços (principal) & HTA 1 (1.1--1.3) e 3.1 (redesenho) & entrada única com busca, ``Mais procurados'' e rota por assunto. \\\hline
    Tela 01a & Validação: campo vazio & HTA 1 & impede busca vazia, destaca a borda e pede que a pessoa escreva o que precisa fazer. \\\hline
    Tela 01b & Validação: texto curto & HTA 1 & ao digitar menos de 3 letras, mantém o texto e pede ao menos 3 caracteres. \\\hline
    Tela 01c & Carregando busca & HTA 1 & indicador de processamento ``Buscando serviços...'' após acionar a pesquisa. \\\hline
    Tela 01d & Busca fora do ar & HTA 1 & avisa a indisponibilidade com ``Tentar de novo'', mantendo ``Mais procurados'' clicável abaixo. \\\hline
    Tela 02 & Resultados da busca & HTA 3.1 & lista os serviços encontrados com nome da tarefa e breve explicação do imposto. \\\hline
    Tela 02a & Vazio: sem resultados & HTA 3.1 & para termo sem retorno, sugere palavras simples, nova busca e atalho para ``Imóveis''. \\\hline
    Tela 03 & Serviços de imóveis & HTA 1.2--1.3 e 3.1 & lista as ações ligadas a imóveis (emitir guia, IPTU, contestar valor, isenção, devolução) em linguagem simples. \\\hline
    Tela 04 & Entrar com gov.br & HTA 2.1 & explica o login único sem senha separada da Receita e o retorno direto ao serviço escolhido. \\\hline
    Tela 04a & Carregando login & HTA 2.1 & exibe ``Entrando com gov.br...'' com botão ``Cancelar login'' visível. \\\hline
    Tela 04b & Erro no login gov.br & HTA 2.1 & informa cancelamento ou demora, tranquiliza que nenhum dado foi enviado e oferece tentar de novo ou voltar. \\\hline
    Tela 05 & Antes de começar & HTA 2 (A8) & apresenta os documentos necessários antes da digitação e confirma os dados trazidos do gov.br. \\\hline
    Tela 06 & Formulário com ITBI escolhido & Término & abre ``Pedir a guia do ITBI'' com \textit{Assunto}, \textit{Tipo de atendimento} e dados cadastrais pré-preenchidos. \\\hline
    Tela 06a & Formulário fora do ar & Término & informa falha na abertura do Atendimento Virtual com opção de tentar novamente, voltar ou ligar 156. \\\hline
    Tela 06b & Validação no formulário & Complementar & demonstra resumo de erros no topo e mensagens específicas nos campos de e-mail, matrícula e descrição. \\\hline
    Tela 06c & Pedido registrado & Complementar & exibe número de protocolo, prazo de envio da guia por e-mail e atalho para consultar o pedido. \\\hline
  \end{tabularx}
  \fonte{\autores}
\end{quadro}\FloatBarrier

% =====================================================================
\chapter{Wireframe da tela principal}\label{cap:wireframe}
% =====================================================================

Haja vista o wireframe de baixa fidelidade (Figura~\ref{fig:wireframe}), que desenvolvemos a partir do esboço, o objetivo nessa fase foi testar a
hierarquia da informação e a sequência de leitura sem nos distrairmos com a ui. No
formato de celular, toda a página foi organizada em coluna única, pensando no
refluxo para telas estreitas sem rolagem horizontal (achado A12, RU03 e RA04),
de modo que a ordem visual coincida com a ordem de navegação por teclado
(achado A11, RA05, RA06 e RA12). A estrutura ficou, assim, dividida em seis blocos:

\begin{enumerate}
  \item Cabeçalho único: identifica o Governo do DF e Receita do DF e o botão de ``Entrar'' e os controles de navegação na mesma posição em todo o fluxo.
  \item Título e instrução com exemplo: traz ``Serviços da Receita do DF'', buscando orientar já na primeira dobra sem depender do jargão interno.
  \item Busca como entrada principal: rótulo sempre visível, dica de palavras cotidianas e o botão ``Buscar'' dimensionado com 40~px de altura para facilitar o toque em tela pequena.
  \item Lista ``Mais procurados'': quatro tarefas comuns nomeadas pelo verbo da ação e acompanhadas da explicação da sigla no subtítulo, destacando ``Emitir guia do ITBI -- Imposto sobre a compra de imóvel (ITBI)''.
  \item Rota alternativa por assunto: seção ``Não achou? Veja os serviços por assunto'' (Imóveis, Veículos, Empresas e Certidões), criada justamente porque em S1 a S4 nenhum participante recorreu à barra de pesquisa de primeira, preferindo explorar categorias (A1, A7 e RU06).
  \item Apoio e rodapé: aviso com telefone 156 e atendimento presencial visível sem sair da página, de mdo a acolher quem precisa de auxílio humano, como apareceu nas sessões S1, S3 e S4 (RU02).
\end{enumerate}

Na adaptação do wireframe para largura de computador, preservamos a
mesma prioridade de leitura, sendo que o cabeçalho passa a mostrar diretamente
os atalhos de ``Acessibilidade'' e ``Alto contraste'', enquanto ``Mais
procurados'' e ``Serviços por assunto'' ficam dispostos lado a lado em duas
colunas logo abaixo da busca.

\fig[0.55]{wireframe}{wireframe_tela_principal.png}{Wireframe da tela principal (Tela 01), largura de celular}{\autores}

% =====================================================================
\chapter{Design visual e protótipo}\label{cap:visual}
% =====================================================================

\section{Componentes}\label{sec:tokens}

Ainda que o serviço pertença ao Governo do Distrito Federal, em que seguir o
Padrão Digital de Governo \cite{dsgov} não é obrigatório, decidimos adotar os
tokens e componentes do DSGov em todo o protótipo. Dois pontos da E1 influenciaram
essa escolha, dado que o padrão resolve de uma vez a fragmentação entre os
três sites atuais da SEEC e as falhas de contraste medidas em A2, A4 e A10.
E, ainda, aproveita padrões que parte dos cidadãos já reconhece de outros
serviços públicos digitais, algo que ocorreu em S2, quando
a participante comparou o portal do GDF ao federal e comentou que neste último
as funções eram bem mais óbvias de achar.

No Figma, cadastramos 29 variáveis semânticas na coleção \textit{DSGov tokens}, além de um modo sem cor usado para o wireframe. Os blocos recorrentes
foram montados como componentes reutilizáveis para manter consistência entre
todas as telas. O Quadro~\ref{qua:tokens-dsgov} sintetiza os tokens e
componentes empregados no fluxo, junto das razões de contraste calculadas pela
fórmula da WCAG 2.2 \cite{wcag22}.

\begin{quadro}[!htbp]
  \caption{Tokens e componentes do DSGov}\label{qua:tokens-dsgov}
  \centering\footnotesize
  \begin{tabularx}{\textwidth}{|P{4.7cm}|P{2.1cm}|Y|}\hline
    Token / componente & Telas & Utilização no fluxo \\\hline
    \texttt{cor/acao-principal} (\texttt{\#1351B4}) e \texttt{acao-principal-escura} (\texttt{\#0C326F}) & 01 a 06c & Botões primários, links de retorno e ações de destaque (7,33:1 e 12,34:1 sobre \texttt{\#FFFFFF}). \\\hline
    \texttt{cor/texto-principal} (\texttt{\#333333}) e \texttt{texto-apoio} (\texttt{\#555555}) & 01 a 06c & Títulos, rótulos, instruções e subtítulos (12,63:1 e 7,46:1 sobre \texttt{\#FFFFFF}). \\\hline
    \texttt{cor/fundo} (\texttt{\#FFFFFF}), \texttt{fundo-alternativo} (\texttt{\#F8F8F8}), \texttt{borda-campo} (\texttt{\#888888}) e \texttt{divisor} (\texttt{\#CCCCCC}) & 01 a 06c & Fundo das páginas, destaque suave, bordas de campos (3,54:1 sobre \texttt{\#FFFFFF}, acima de 3:1) e linhas divisórias decorativas. \\\hline
    \texttt{cor/erro} (\texttt{\#E52207}) e \texttt{erro-fundo} (\texttt{\#FDE0DB}) & 01a, 01b, 01d, 04b, 06a, 06b & Bordas, ícones e texto de erro sobre branco (4,60:1). Nos cartões com \texttt{erro-fundo}, o texto usa \texttt{\#333333} (10,14:1) e o ícone/borda mantém 3,69:1. \\\hline
    \texttt{cor/sucesso} (\texttt{\#168821} / \texttt{\#E3F5E1}), \texttt{informacao} (\texttt{\#155BCB} / \texttt{\#D4E5FF}) e \texttt{alerta} (\texttt{\#FFCD07} / \texttt{\#FFF5C2}) & 01, 02a, 04, 05, 06, 06c & Caixas de orientação, aviso de dados do gov.br, estado vazio e protocolo. Textos internos em \texttt{\#333333} garantem contraste superior a 10:1 (ex.: 11,49:1 sobre \texttt{\#FFF5C2}). \\\hline
    \texttt{cor/foco} (\texttt{\#C2850C}), \texttt{rodape} (\texttt{\#071D41}) e \texttt{texto-sobre-escuro} (\texttt{\#FFFFFF}) & 01 a 06c & Contorno de foco de teclado (3,16:1 sobre \texttt{\#FFFFFF}, atendendo 1.4.11) e rodapé institucional (16,65:1). \\\hline
    \texttt{espaco} (4, 8, 16, 24, 32 e 40~px), \texttt{raio} (4, 8 e 100~px) e \texttt{altura/alvo} (40~px) & 01 a 06c & Separação entre blocos e campos, arredondamento padronizado e altura mínima de 40~px para alvos clicáveis, superando o mínimo de 24~px de RA07. \\\hline
    Cabeçalho (\texttt{br-header}) e Rodapé (\texttt{br-footer}) & 01 a 06c & Identidade única da Receita do DF, login/usuário logado e canais de ajuda (156 e presencial). \\\hline
    Campo (\texttt{br-input}) e Área de texto (\texttt{br-textarea}) & 01 a 02, 06, 06b & Busca de serviços, dados cadastrais e descrição do pedido com limite de caracteres e instrução de saída por Tab. \\\hline
    Botões (\texttt{br-button}) e Lista (\texttt{br-list} / \texttt{br-item}) & 01 a 06c & Ações principais e secundárias; listas clicáveis com título da tarefa, explicação da sigla e seta indicativa. \\\hline
    Mensagens (\texttt{br-message}), Seletores (\texttt{br-radio} / \texttt{br-checkbox}) e Carregamento (\texttt{br-loading}) & 01 a 06c & Alertas de erro, validação, vazio e indisponibilidade; opções do formulário e feedback de processamento. \\\hline
  \end{tabularx}
  \fonte{\autores, com base em \citeonline{dsgov}.}
\end{quadro}\FloatBarrier

Em relação ao padrão federal original, fizemos duas mudanças, as quais sao justificadas pelos achados da E1:

\begin{itemize}
  \item No lugar da marca ``gov.br'', o cabeçalho exibe ``GDF -- Governo do Distrito Federal / Receita do DF / Secretaria de Economia do DF''. A mudança reflete a jurisdição do serviço, mantendo a mesma anatomia e previsibilidade do DSGov do início ao fim do percurso.
  \item No serviço real, a área ``Descrição da Solicitação'' abre um editor de texto completo com mais de uma dezena de botões sem rótulo que ainda prende o foco do teclado (A9 e A11). No protótipo, reduzimos a barra de formatação e adicionamos a instrução explícita ``Até 4.000 caracteres. Use Tab para sair do editor.'', de modo a cumprir RA09, RA10 e RA11.
\end{itemize}

Vale observar, ainda, que as ações secundárias do Atendimento Virtual real
(``Consultar pedido'', ``Reabrir pedido'', ``Complementar pedido'', ``Avaliar
atendimento'' e ``Passo a passo''), que antes distraiam o usuário no topo da tela com a
ação principal (achado A8), foram deslocadas para a base do formulário na Tela
06, de modo a preservar as funcionalidades do órgão sem encher demais a entrada com várias opções.

\section{Telas}\label{sec:estados}

Todas as telas foram projetadas primeiro na largura de celular, levando em conta que o smartphone é o principal meio de acesso à
internet no Brasil \cite{ticdomicilios} e que S4 relatou excesso de botões e
informação ao tentar o fluxo pelo telefone. Embora o quadro no Figma use 360~px
(atendendo com folga ao recorte de até 412~px de RU03), o layout de coluna
única sem tabelas largas foi desenhado para suportar o refluxo até 320~px sem
rolagem horizontal (RA04). A tela principal (Figura~\ref{fig:principal-desktop})
e o restante do fluxo navegável também possuem versão para computador, de modo a distruibuir ``Mais procurados'' e ``Serviços por assunto'' lado a
lado.

\fig[0.9]{principal-desktop}{tela_principal_desktop.png}{Tela principal (Tela 01) na versão web}{\autores}

No conteúdo, evitamos o uso de texto fictício genérico dado que as telas usam
exemplos plausíveis para o contexto do DF (solicitante fictícia Maria da Silva
Souza, protocolo \texttt{2026.0012345}, matrícula \texttt{123456}). Além disso, aplicamos as diretrizes da Política Nacional de Linguagem Simples
\cite{linguagemsimples2025}, usando frases curtas, voz ativa e termos
cotidianos ao inves do jargão da carta de serviços atual (como trocar
``Transmissão Onerosa (Lançamento de ITBI)'' por ``Emitir guia do ITBI (Imposto sobre a compra de imóvel)'').

Quanto aos estados de interação e falha, os cenários previstos
foram representados nas telas do protótipo de modo a evitar deixar o
usuário sem saber como proceder quando algo dá errado:

\begin{itemize}
  \item Validação de campo obrigatório vazio (Tela 01a e Tela 06b): ao acionar ``Buscar'' sem digitar nada, o campo fica borda vermelha, ícone de alerta e a mensagem ``Escreva o que você precisa fazer para buscar.''. De forma complementar, na Tela 06b o resumo aponta os três campos pendentes e o campo de descrição vazio mostra instrução com exemplo.
  \item Validação de formato inválido (Tela 01b e Tela 06b): se a pessoa digitar menos de três letras na busca, o sistema mantém o que foi escrito e orienta ``Escreva pelo menos 3 letras. Exemplo: ITBI.''. Na Tela 06b, e-mails diferentes e matrícula preenchida com letras ou traço mostram uma correção específica proximo ao campo.
  \item Carregando (Telas 01c e 04a): durante a consulta de serviços aparece ``Buscando serviços...'' e, na autenticação, ``Entrando com gov.br...'' acompanhado do botão ``Cancelar login'' para que a pessoa não fique presa em caso de lentidão.
  \item Estado vazio (Tela 02a): quando a pesquisa não encontra correspondência, a tela informa ``Nenhum serviço encontrado'', sugere termos simples (``guia, imposto ou imóvel''), oferece ``Fazer nova busca'' e já entrega um atalho direto para a categoria ``Imóveis''.
  \item Erro de autenticação (Tela 04b): se o login no gov.br for cancelado ou expirar, a mensagem explica o ocorrido (``Você não concluiu o login''), assegura que nenhum dado foi enviado e oferece ``Tentar de novo'' ou ``Voltar para o início''.
  \item Serviço fora do ar (Telas 01d e 06a): se a busca falhar (01d), o aviso ``A busca está fora do ar'' traz o botão ``Tentar de novo'' e mantém a lista ``Mais procurados'' logo abaixo como rota alternativa. Se o próprio Atendimento Virtual estiver indisponível (06a), a tela informa o problema e oferece tentar de novo, voltar ao início ou ligar no 156.
\end{itemize}

\ifcodigo
\chapter{Validação na tela em código}\label{cap:validacao}
\chapter{Auditoria automática de acessibilidade}\label{cap:auditoria}
\fi

% =====================================================================
\chapter{Decisões de design}\label{cap:decisoes}
% =====================================================================

As decisões principais do redesenho foram documentadas, de forma gera, seguindo o formato QOC
(\textit{Questions, Options and Criteria}), proposto por
\citeonline{maclean1991}. Esse método ajuda a deixar explícito não só o que entrou
no protótipo, mas também qual alternativa foi descartada e quais requisitos e
evidências da E1 pesaram na escolha. Assim, D1 registra a escolha da alternativa da tela
principal a partir dos esboços em sala e D2 trata das condições de acesso
precário (no caso, a persona 6, Edilene). Já D3, D4 e D5 abrangem o vocabulário em linguagem
simples, a unificação visual com contraste acessível e a reorganização dos
pré-requisitos dada a simplificação do formulário.

% D1 -------------------------------------------------------------------
\begin{quadro}[!htbp]
  \caption{Decisão D1, alternativa adotada para a tela principal}\label{qua:d1}
  \centering\small
  \antesdepois{fig_home.png}{E1, Figura 1 (menu Serviços em cascata)}{tela_01_mobile.png}{Tela 01 (Busca de serviços)}
  \qoc{Como organizar a tela inicial para que tanto quem prefere pesquisar quanto quem prefere navegar ache o serviço de ITBI sem precisar decifrar menus administrativos?}%
      {Alternativa A: página de entrada única com busca em destaque no topo, seguida por lista curta de ``Mais procurados'' nomeados pela tarefa e acesso secundário por assunto.}%
      {Alternativa B (esboço descartado): entrada estruturada por menu lateral/superior dividido por perfil de público (\textit{Atendimento Cidadão}, \textit{Atendimento Empresa}) e categorias hierárquicas (\textit{Serviços} $\rightarrow$ \textit{Imóveis} $\rightarrow$ \textit{ITBI}), sem destaque para a busca.}%
      {RU01, RU02, RU04, RU06, RA01, RA05, RA12}%
      {Achados A1, A5, A7 e A11; sessões S1 a S5; personas 1, 3, 4, 5 e 7.}%
      {Reunir busca visível, tarefas frequentes na primeira dobra e categorias simples reduz o número de escolhas simultâneas (Lei de Hick-Hyman) \cite{hick1952}, atende aos dois estilos de uso vistos nas sessões e elimina o menu em cascata que só abria por \textit{hover}, permitindo uma ordem de foco previsível por teclado.}
  \fonte{\autores}
\end{quadro}\FloatBarrier

% D2 -------------------------------------------------------------------
\begin{quadro}[!htbp]
  \caption{Decisão D2, condições de acesso precário}\label{qua:d2}
  \centering\small
  \antesdepois{fig_cat.png}{E1, Figura 2 (catálogo em três colunas e múltiplos carregamentos)}{tela_01d_mobile.png}{Tela 01d (coluna única e rota alternativa)}
  \qoc{Como viabilizar que uma pessoa em celular de entrada, com tela pequena e internet móvel limitada ou instável, chegue ao formulário sem gastar franquia à toa nem travar em falhas pontuais?}%
      {Layout em coluna única projetado para refluxo até 320~px sem rolagem horizontal, remoção de banners pesados, alvos clicáveis de 40~px de altura (\texttt{altura/alvo}), menos carregamentos em domínio único e tela ``Fora do ar'' (01d) que preserva ``Mais procurados'' clicável mesmo quando a busca cai.}%
      {Manter os três subdomínios, os quatro carregamentos intermediários, o banner institucional da NFS-e e o catálogo em três colunas apenas ajustados por redimensionamento simples.}%
      {RU03, RA04, RA07}%
      {Achados A2, A5 e A12; sessão S4; personas 2 (Júlia) e 6 (Edilene, acesso precário).}%
      {Cortar recargas de página e elementos pesados economiza dados móveis e tempo de carregamento; alvos de 40~px reduzem erros de toque em telas pequenas, na linha da Lei de Fitts \cite{fitts1954} e acima do piso de 24~px da WCAG 2.2 \cite{wcag22}; e manter os atalhos na Tela 01d impede que uma instabilidade na busca bloqueie a jornada como um todo.}
  \fonte{\autores}
\end{quadro}\FloatBarrier

% D3 -------------------------------------------------------------------
\begin{quadro}[!htbp]
  \caption{Decisão D3, vocabulário mais acessível}\label{qua:d3}
  \centering\small
  \antesdepois{a3_catalogo.png}{E1, Figura 9 (catálogo mais jurídico)}{tela_03_mobile.png}{Tela 03 (Linguagem simples)}
  \qoc{Como nomear os serviços ligados a imóveis para que o cidadão reconheça a emissão da guia sem dominar jargão jurídico nem confundir telas diferentes chamadas apenas ``ITBI''?}%
      {Nomear cada opção pela ação concreta em linguagem simples (``Emitir guia do ITBI'', ``Contestar o valor do ITBI'', ``Pedir isenção do ITBI''), explicar a sigla no próprio subtítulo (``Imposto sobre a compra de imóvel'') usando texto real em vez de cartões em imagem e manter o termo técnico apenas como apoio (``Também chamado de lançamento de ITBI'').}%
      {Manter os rótulos jurídicos atuais (``Transmissão Onerosa -- Lançamento de ITBI'', ``Efetuar Impugnação'') e tentar compensar com um glossário ou ícone de ajuda à parte.}%
      {RU01, RU02, RU05, RU06, RA08, RA10}%
      {Achados A3, A6 e A10; sessões S2, S4 e S5; personas 1, 2, 4, 5 e 7.}%
      {Substituir termos internos pelo verbo da tarefa auxilia no reconhecimento ao invés da recordação \cite{norman2013}, segue a Lei de Linguagem Simples \cite{linguagemsimples2025} e comabte a causa comum de seis achados da E1, além de eliminar os cartões em imagem sem texto alternativo que bloqueavam o leitor de tela.}
  \fonte{\autores}
\end{quadro}\FloatBarrier

% D4 -------------------------------------------------------------------
\begin{quadro}[!htbp]
  \caption{Decisão D4, unificação visual sob o DSGov e correção de contraste}\label{qua:d4}
  \centering\small
  \antesdepois{c_cabecalhos.png}{E1, Figura 8 (barra de acessibilidade ausente)}{tela_01_desktop.png}{Telas 01 a 06 (contraste nível AA)}
  \qoc{Como tirar a sensação de mudança de site no meio do caminho e tornar textos, bordas, estados ativos e foco legíveis para quem tem baixa visão ou tela com pouco brilho?}%
      {Adotar cabeçalho único e tokens de cor, tipografia e foco do DSGov em todas as telas do fluxo, mantendo os controles de login e acessibilidade na mesma posição e garantindo contraste mínimo nível AA (pelo menos 4,5:1 em textos e 3:1 em bordas e foco).}%
      {Preservar as três identidades visuais separadas (portal institucional, Receita e Atendimento Virtual), trocando apenas a cor dos botões que reprovaram na medição de contraste.}%
      {RU07, RU08, RA02, RA03, RA06, RA08}%
      {Achados A2, A4 e A10; sessões S1, S2 e S3; personas 1, 3 e 5.}%
      {Unificar o cabeçalho e o sistema visual sob o DSGov \cite{dsgov} auxilia quanto ao o senso de localização do usuário desde o início até o fim da sua navegabilidade, impede que a acessibilidade não exista na etapa do formulário e corrige os contrastes críticos de 1,35:1 a 2,61:1 encontrados na E1.}
  \fonte{\autores}
\end{quadro}\FloatBarrier

% D5 -------------------------------------------------------------------
\begin{quadro}[!htbp]
  \caption{Decisão D5, pré-requisitos e simplificação do formulário}\label{qua:d5}
  \centering\small
  \antesdepois{a9_form.png}{E1, Figura 15 (seletores vazios)}{tela_06_mobile.png}{Telas 04, 05 e 06 (formulário pré-selecionado)}
  \qoc{Como avisar login e documentos antes do preenchimento e reduzir tanto a carga cognitiva quanto as barreiras de teclado ao abrir o formulário de solicitação?}%
      {Inserir a entrada única via gov.br (Tela 04) e a tela ``Antes de começar'' com a lista de documentos (Tela 05), abrindo a Tela 06 com \textit{Assunto}, \textit{Tipo de atendimento} e dados cadastrais já preenchidos, sinalizando apenas campos ``(opcional)'' em texto claro e simplificando a descrição com saída por Tab e validação explícita (Tela 06b).}%
      {Enviar a pessoa direto ao formulário legado com seletores dependentes em ``Selecione'', asterisco vermelho claro sem associação programática e editor de texto completo, mantendo o ``Passo a Passo'' escondido entre seis botões iguais no topo.}%
      {RU04, RU05, RU07, RU08, RA01, RA09, RA10, RA11}%
      {Achados A8, A9 e A11; sessões S1 e S4; inspeção com NVDA e teclado; personas 1, 3, 4, 5 e 6.}%
      {Informar os pré-requisitos antes da digitação pode evitar a desistência tardia e idas desnecessárias ao atendimento presencial; e reduzir a descrição remove os botões sem rótulo e a armadilha de teclado identificada na E1.}
  \fonte{\autores}
\end{quadro}\FloatBarrier

% =====================================================================
\chapter{Matriz de rastreabilidade}\label{cap:matriz}
% =====================================================================

A matriz de rastreabilidade atualizada para a entrega 2 foi entregue, no Moodle (\texttt{UCEngUs\_\allowbreak matriz\_\allowbreak GrupoB.xlsx}).

% =====================================================================
\chapter{Declaração de contribuição individual}\label{cap:contribuicao}
% =====================================================================

\begin{quadro}[!htbp]
  \caption{Contribuição individual}\label{qua:contribuicao}
  \centering\small
  \begin{tabularx}{\textwidth}{|P{4.2cm}|Y|}\hline
    Integrante & Atividades realizadas \\\hline
    Ícaro Souza & Elaboração da matriz de rastreabilidade e apoio na revisão geral. \\\hline
    Karoline Costa & Escrita do memorial e construção do protótipo no Figma. \\\hline
    Leonardo Brito & Escrita do memorial e construção do protótipo no Figma. \\\hline
    Luiz Gustavo & Escrita do memorial e construção do protótipo no Figma. \\\hline
  \end{tabularx}
  \fonte{\autores}
\end{quadro}\FloatBarrier

% =====================================================================
\chapter{Declaração de uso de inteligência artificial}\label{cap:ia}
% =====================================================================

\begin{quadro}[!htbp]
  \caption{Uso de inteligência artificial}\label{qua:ia}
  \centering\small
  \begin{tabularx}{\textwidth}{|P{3.6cm}|Y|P{2.8cm}|}\hline
    Ferramenta & Forma de uso & Etapa \\\hline
    Claude (Anthropic) & Apoio na geração inicial da estrutura \LaTeX{} (capítulos e quadros) deste memorial, revisada pelos autores. & Diagramação \\\hline
  \end{tabularx}
\end{quadro}\FloatBarrier

\postextual

\bibliography{referencias}

\begin{anexosenv}
\partanexos

\chapter{Esboços}\label{an:esbocos}

\fig[0.7]{esboco-icaro}{esboco_icaro.jpg}{Esboço do Ícaro Souza}{\autores}
\fig[0.7]{esboco-karoline}{esboco_karoline.jpg}{Esboço da Karoline Costa}{\autores}
\fig[0.7]{esboco-leonardo}{esboco_leonardo.jpg}{Esboço do Leonardo Brito}{\autores}
\fig[0.7]{esboco-luiz}{esboco_luiz.jpg}{Esboço do Luiz Gustavo}{\autores}

\end{anexosenv}

\end{document}
